% Original analytic teaching diagram, 2026-11-09.
% BFGS 逆更新的分解：s=(24,14)，y=(28,6)，z=(12,24)，旧矩阵示意取单位矩阵；比例 a=s^Tz/s^Ty=0.8254。
\definecolor{VB}{HTML}{7188A3}\definecolor{VBF}{HTML}{E5EAF0}
\definecolor{VG}{HTML}{829C8B}\definecolor{VGF}{HTML}{E8EEE9}
\definecolor{VY}{HTML}{B4A16F}\definecolor{VYF}{HTML}{F1ECDD}
\definecolor{VR}{HTML}{A7787D}\definecolor{VRF}{HTML}{F0E5E6}
\definecolor{VP}{HTML}{9588A5}\definecolor{VPF}{HTML}{ECE8F0}
\definecolor{VI}{HTML}{3B4252}\definecolor{VM}{HTML}{6B7280}\definecolor{VL}{HTML}{4C4D4F}
\begin{tikzpicture}[x=1mm,y=1mm,>=Stealth,font=\figurefont\fontsize{8}{10}\selectfont,
lab/.style={text=VM,font=\figurefont\fontsize{7.5}{9}\selectfont},
vec/.style={->,line width=1.1pt},aux/.style={VL,line width=.6pt,dash pattern=on 2pt off 1.5pt}]
\begin{scope}
\draw[VM,line width=.6pt] (11.2,-19.2)--(-16.8,28.8) node[pos=.04,right,lab] {$\bm{s}$的正交补};
\draw[VM,line width=.6pt] (-4.2,-.9)--(35,7.5);
\draw[aux] (23.11,4.95)--(12,24)--(-11.11,19.05);
\draw[vec,VG] (0,0)--(23.11,4.95) node[midway,below,text=VI] {$a\bm{y}$};
\draw[vec,VP] (0,0)--(-11.11,19.05) node[left,text=VI] {$\bm{V}\bm{z}$};
\draw[vec,VI] (0,0)--(12,24) node[above,text=VI] {$\bm{z}$};
\draw[vec,VB,line width=.9pt] (0,0)--(28,6) node[right,text=VI] {$\bm{y}$};
\fill[VL] (0,0) circle (.5);
\node[text=VI] at (8,-23.5) {(a)};
\end{scope}
\draw[->,VL,line width=.75pt] (37,14)--(48,14) node[midway,above,text=VI] {$\bm{W}^{+}$};
\begin{scope}[shift={(61,-2)}]
\draw[VM,line width=.6pt] (5.4,-25.2+8)--(-7,32.67) node[pos=.03,right,lab] {$\bm{y}$的正交补};
\draw[VM,line width=.6pt] (-4.8,-2.8)--(30,17.5);
\draw[aux] (19.81,11.56)--(14.95,34.25)--(-4.86,22.69);
\draw[vec,VG] (0,0)--(19.81,11.56) node[midway,below right,xshift=-1mm,text=VI] {$a\bm{s}$};
\draw[vec,VP] (0,0)--(-4.86,22.69) node[left,text=VI] {$\bm{V}^\top\bm{W}\bm{V}\bm{z}$};
\draw[vec,VI] (0,0)--(14.95,34.25) node[right,text=VI] {$\bm{W}^{+}\bm{z}$};
\draw[vec,VB,line width=.9pt] (0,0)--(24,14) node[right,text=VI] {$\bm{s}$};
\fill[VL] (0,0) circle (.5);
\node[text=VI] at (8,-21.5) {(b)};
\end{scope}
\end{tikzpicture}
